\subsection{用阶梯函数逼近连续函数}

命题 19. --- 设 X 是一个局部紧空间，\(\Phi\) 是 X 的子集所成的一个集环，它包含 X 的紧子集全体所成之集。对每个具有紧支集 K 的、X 到 Banach 空间 F 中的连续映射 f（相应地，X 上的每个连续实值函数 \(f \geqslant 0\)），存在函数的一个序列 \((\mathbf{g}_n)\)（属于 \(\mathcal{E}_{\mathrm{F}}(\Phi)\)，其支集含于 K）（相应地，存在函数的一个序列 \((g_n)\)（属于 \(\mathcal{E}(\Phi)\)，使得 \(0 \leqslant g_n \leqslant f\) 对每个 n 成立）），它一致收敛于 f（相应地 \(f\)）。

由于 \(\mathbf{f}\) 在 \(K\) 上一致连续，\(K\) 可被有限个紧集 \(\mathbf{M}_i\) （ \(1 \leqslant i \leqslant m\) ）所覆盖，使得 \(\mathbf{f}\) 在每个 \(\mathbf{M}_i\) 上的振幅 \(\leqslant 1/n\)。由于 \(\mathbf{M}_i\) 与 \(K\) 属于 \(\Phi\)，存在把 \(K\) 分成集合 \(\mathbf{N}_j \in \Phi\) 的一个分割，使得每个集合 \(\mathbf{M}_i \cap K\) 是若干个 \(\mathbf{N}_j\) 的并（No.~9, 引理）。设 \(\mathbf{a}_j\) 是 \(F\) 的一个元素，使得 \(|\mathbf{f}(x) - \mathbf{a}_j| \leqslant 1/n\) 在 \(\mathbf{N}_j\) 上成立。令 \(\mathbf{g}_n = \sum_j \mathbf{a}_j \varphi_{\mathbf{N}_j}\)，我们就有 \(|\mathbf{f} - \mathbf{g}_n| \leqslant 1/n\)，由此得到该命题在这一情形的结论。对连续实值函数 \(f\) 可类似论证，只需取 \(a_j = \inf_{x \in \mathbb{N}_j} f(x)\) 与 \(g_n = \sum_j a_j \varphi_{\mathbf{N}_j}\)。

推论 1. --- 设 \(\mu\) 是 X 上的一个测度；空间 \(\mathcal{E}_{\mathrm{F}}(\Phi)\) 在每个空间 \(\mathcal{L}_{\mathrm{F}}^{p}\) （ \(1 \leqslant p < +\infty\)）中稠密。

因为，由命题 19 以及具有紧支集的函数的一致极限的平均收敛准则（§3, No.~3, 命题 4）得出，\(\mathcal{E}_{\mathrm{F}}(\Phi)\) 关于 \(p\) 阶平均收敛拓扑，在具有紧支集的连续函数所成的空间 \(\mathcal{K}_{\mathrm{F}}\) 的闭包中稠密，由此即得该推论。

推论 2. --- 对 X 的每个闭子集 S，每个函数 \(f \in \mathcal{K}(X, S; \mathbf{C})\) 都是线性组合 \(\sum_{i} \lambda_i \varphi_{K_i}\) 的一致极限，其中 \(\lambda_i\) 属于 \(\mathbf{C}\)，而 \(K_i\) 是 \(S\) 的紧子集。

这种线性组合所成之集 A 是一个 C-代数。设 \(\Phi\) 是 X 的使 \(\varphi_{M} \in A\) 的子集 M 所成之集；这样 \(\Phi\) 就是一个集环，其所有元素都是 S 的子集，它包含 S 的紧子集，并且 \(\mathcal{E}_{\mathbf{C}}(\Phi) \subset \mathcal{A}\)。于是只需对局部紧空间 S 与集环 \(\Phi\) 应用命题 19。

推论 3. --- 若 \(\mu\) 与 \(\nu\) 是 X 上的两个测度，使得 \(\mu(K) = \nu(K)\) 对每个紧子集 \(K\)（\(X\) 的）成立，则 \(\mu = \nu\)。

因为，由推论 2 与测度的定义得出，对 X 的每个紧子集 S，\(\mu\) 与 \(\nu\) 在 \(\mathcal{K}(\mathrm{X},\mathrm{S};\mathbf{C})\) 上取相同的值。

